% Propietario conservado: /Users/ruben/Documents/New project/output/REVISION_ARGUMENTAL_APERTURAS_CIERRES_20260915/VII/delivery/MANUSCRIPT_VII_ES_EN_COMPLETE/source_es/manuscrito/30b_variacion_energia_memoria.tex
% Artículo VII. Desarrollo aditivo, revisión 04, 10-09-2026.
% ARQUITECTURA_AUTORAL_PREEXISTENTE: transporte enriquecido, energía
% D_9^* W D_9 y compatibilidad por refinamiento, propietario S01,
% 02_dinamica_tres_hojas.tex, líneas 221--250.
% FORMALIZACION_NUEVA / CERTIFICADO_NUEVO: diferencial explícito de ese
% funcional, representante antihermítico, composición e inversión ponderada.
% Las variaciones pertenecen al espacio de realizaciones unitarias del
% transporte. No se presume que toda variación continua sea una ruta TPK.
% Su realización como corriente material de espín se trata por separado.
% \mathsf W_m es el peso energético; W de la sección anterior es soldadura.

\section{Variational response of memory energy}
\label{vii:vii:sec:variacion-memoria}

The transport of a route identifies the fibers at its endpoints and preserves
the orientation and memory entering that identification. The associated
energy compares the final field with the transported initial field. Its
differential quantifies how that comparison responds to a variation of
transport. This section obtains the response directly from the discrete
functional and develops its composition, covariance and inversion.

\subsection{Discrete functional and space of variations}
\label{vii:vii:subsec:funcional-memoria}

At a finite level \(m\), let \(\mathcal V_m\) be the vertices of the
route publication and \(\mathcal E_m^+\) an orientation of each edge.
Each vertex carries a finite-dimensional Hermitian fiber \(H_v\).
We define
\[
 \mathcal H_m=\bigoplus_{v\in\mathcal V_m}H_v,\qquad
 \mathcal K_m=\bigoplus_{a\in\mathcal E_m^+}H_{t(a)}.
\]
The Hermitian inner product is antilinear in the first variable.
For unitary \(U_a:H_{s(a)}\to H_{t(a)}\), the difference operator is
\begin{equation}
 (D_mf)_a=f_{t(a)}-U_af_{s(a)},\qquad
 E_m(f,U)=\langle D_mf,\mathsf W_mD_mf\rangle,\qquad
 \mathsf W_m=\mathsf W_m^*>0.
 \label{vii:vii:eq:energia-memoria-discreta}
\end{equation}
This is the energy of the term \(D_m^*\mathsf W_mD_m\) in the discrete generator.
The weight \(\mathsf W_m\) may couple different edges. The soldering operator \(W\)
of Section~\ref{vii:vii:sec:pantalla-respuesta} and this energy
weight have different domains and roles.

The transport \(U\) is a publication of the enriched state. The field
\(f\) and the weight \(\mathsf W_m\) remain explicit as arguments
of the functional; each material application preserves the rules that
determine them. Differentiable families of the unitary realizations
of transport are considered,
\[
 \dot U_a=A_aU_a,\qquad A_a^*=-A_a.
\]
The variations used physically form the subspace preserving
the conditions of the corresponding realization. The differential is
calculated first in the unitary tangent space and then restricted
to that subspace.

\begin{proposition}[Transport-response covector]
\label{vii:vii:prop:covector-memoria}
Write
\[
 v_a=U_af_{s(a)},\qquad r=D_mf,\qquad q=\mathsf W_mr.
\]
At fixed field and weight, the energy differential is
\begin{equation}
 \mathfrak j_m(A)=-2\operatorname{Re}
       \sum_{a\in\mathcal E_m^+}\langle q_a,A_av_a\rangle .
 \label{vii:vii:eq:covector-memoria}
\end{equation}
Its representative for the real Hilbert--Schmidt inner product is
\begin{equation}
 \mathcal J_a=v_aq_a^*-q_av_a^*,\qquad
 \mathcal J_a^*=-\mathcal J_a,\qquad
 \mathfrak j_m(A)=
 \sum_a\operatorname{Re}\operatorname{Tr}(\mathcal J_a^*A_a).
 \label{vii:vii:eq:representante-corriente-discreta}
\end{equation}
For simultaneous variations of field, transport and weight we have
\begin{equation}
 \dot E_m=
 2\operatorname{Re}\sum_a
 \left\langle q_a,\dot f_{t(a)}-U_a\dot f_{s(a)}-A_av_a\right\rangle
 +\langle r,\dot{\mathsf W}_mr\rangle .
 \label{vii:vii:eq:variacion-completa-memoria}
\end{equation}
\end{proposition}

\begin{proof}
The first-order update, written in the algebra of dual
numbers with \(\epsilon^2=0\), gives
\[
 r+\epsilon\dot r,\qquad
 \dot r_a=\dot f_{t(a)}-U_a\dot f_{s(a)}-A_av_a.
\]
The coefficient of \(\epsilon\) in the energy is
\[
 \langle\dot r,\mathsf W_mr\rangle+
 \langle r,\mathsf W_m\dot r\rangle+
 \langle r,\dot{\mathsf W}_mr\rangle.
\]
Self-adjointness of the weight yields
\eqref{vii:vii:eq:variacion-completa-memoria}. Fixing field and weight leaves
\eqref{vii:vii:eq:covector-memoria}. Moreover,
\[
 \operatorname{Tr}(\mathcal J_a^*A_a)
 =v_a^*A_aq_a-q_a^*A_av_a.
\]
Anti-Hermiticity of \(A_a\) implies
\(v_a^*A_aq_a=-\overline{q_a^*A_av_a}\), giving
the stated representative.
\end{proof}

The response is thus evaluated: on each edge the field is transported,
its difference from the final field is calculated and the energy weight is applied.
The antisymmetrized product of those two vectors determines
\(\mathcal J_a\in\mathfrak u(H_{t(a)})\). This operator represents,
through the real Hilbert--Schmidt inner product, the response covector
on the unitary variations of the target fiber.

\subsection{Composition and transport of the response}
\label{vii:vii:subsec:composicion-respuesta}

\begin{proposition}[Composition rule along a route]
\label{vii:vii:prop:composicion-respuesta}
For a route formed by \(N\) transports, let
\(U_\gamma=U_N\cdots U_1\) and
\(B_k=U_N\cdots U_{k+1}\), with \(B_N=I\).
If \(\dot U_k=A_kU_k\), then
\begin{equation}
 \dot U_\gamma U_\gamma^{-1}
 =\sum_{k=1}^N B_kA_kB_k^{-1}.
 \label{vii:vii:eq:variacion-transporte-compuesto}
\end{equation}
Thus a representative \(\mathcal J_\gamma\) in the final fibre
induces in factor \(k\) the contribution
\(B_k^*\mathcal J_\gamma B_k\). When a factor belongs to several
routes in the functional, its contributions are summed with the incidences
used by that functional.
\end{proposition}

\begin{proof}
The product rule gives
\[
 \dot U_\gamma=\sum_k
 U_N\cdots U_{k+1}A_kU_k\cdots U_1.
\]
Multiplication by \(U_\gamma^{-1}\) eliminates the factors to the right
of \(A_k\) and yields \eqref{vii:vii:eq:variacion-transporte-compuesto}.
By unitarity and cyclicity of the trace,
\[
 \operatorname{Re}\operatorname{Tr}
   (\mathcal J_\gamma^*B_kA_kB_k^{-1})
 =\operatorname{Re}\operatorname{Tr}
   ((B_k^*\mathcal J_\gamma B_k)^*A_k).
\]
Linearity of the differential proves the sum of contributions.
\end{proof}

This formula preserves the position of each operation within the route.
Transport following an incidence determines how
its contribution is represented in the final fiber; the local response is recovered
by transporting it to the fiber of that incidence.

\begin{proposition}[Covariance under unitary frame changes]
\label{vii:vii:prop:covariancia-corriente-memoria}
Let \(g_v\) be unitary frame changes independent of the variation
parameter, and \(G_{\mathcal E}=\bigoplus_a g_{t(a)}\).
Under the transformations
\[
 f'_v=g_vf_v,\quad U'_a=g_{t(a)}U_ag_{s(a)}^{-1},\quad
 \mathsf W'_m=G_{\mathcal E}\mathsf W_mG_{\mathcal E}^{-1},\quad
 A'_a=g_{t(a)}A_ag_{t(a)}^{-1},
\]
we have
\begin{equation}
 \mathcal J'_a=g_{t(a)}\mathcal J_ag_{t(a)}^{-1},
 \qquad \mathfrak j'_m(A')=\mathfrak j_m(A).
 \label{vii:vii:eq:covariancia-respuesta}
\end{equation}
\end{proposition}

\begin{proof}
We obtain \(r'=G_{\mathcal E}r\), \(q'=G_{\mathcal E}q\) and
\(v'_a=g_{t(a)}v_a\). Substitution into
\eqref{vii:vii:eq:representante-corriente-discreta} proves the first
identity. Invariance of the Hermitian inner product, or cyclicity of
the trace, proves the second.
\end{proof}

When frame changes depend on the parameter, their derivatives
also act on the field and weight. The corresponding covariant
expression is then the complete differential
\eqref{vii:vii:eq:variacion-completa-memoria}.

\subsection{Oriented inversion and transport of the weight}
\label{vii:vii:subsec:inversion-respuesta}

Inversion can be seen explicitly in the case of a local energy
per edge, \(\mathsf W_m=\bigoplus_a W_a\).
For \(a:s\to t\), its inverse carries
\[
 U_{\bar a}=U_a^{-1},\qquad
 r_{\bar a}=-U_a^{-1}r_a,\qquad
 W_{\bar a}=U_a^{-1}W_aU_a.
\]
These rules preserve the edge energy exactly.

\begin{proposition}[Differential compatibility with inversion]
\label{vii:vii:prop:inversion-respuesta}
If \(W_a\) and the fields at both endpoints remain fixed, then
\[
 A_{\bar a}=-U_a^{-1}A_aU_a,\qquad
 \dot W_{\bar a}=U_a^{-1}[W_a,A_a]U_a.
\]
The complete variation of the reversed edge's energy coincides with
the original. In particular, if \([W_a,A_a]=0\), the partial transport
covector satisfies
\begin{equation}
 \mathfrak j_{\bar a}(U_a^{-1}A_aU_a)
 =-\mathfrak j_a(A_a).
 \label{vii:vii:eq:imparidad-transporte}
\end{equation}
\end{proposition}

\begin{proof}
The derivative of \(U_a^{-1}\) is \(-U_a^{-1}A_a\), from which
the two formulas result. Writing \(v_a=U_af_s\) and \(f_t=r_a+v_a\),
the part due exclusively to reversed transport equals
\[
 \mathfrak j_{\bar a}(A_{\bar a})
 =-2\operatorname{Re}\langle r_a,W_aA_af_t\rangle
 =\mathfrak j_a(A_a)
   -2\operatorname{Re}\langle r_a,W_aA_ar_a\rangle.
\]
Variation of the weight contributes
\[
 \langle r_{\bar a},\dot W_{\bar a}r_{\bar a}\rangle
 =\langle r_a,[W_a,A_a]r_a\rangle
 =2\operatorname{Re}\langle r_a,W_aA_ar_a\rangle.
\]
The two terms compensate each other. If the commutator is zero, that contribution
disappears; linearity in \(A_a\) gives
\eqref{vii:vii:eq:imparidad-transporte}.
\end{proof}

Energy is summed over one orientation per edge. A sum over both
orientations carries the factor \(1/2\). This convention preserves the same
functional and hence the same response.

\subsection{Compatibility with refinement}
\label{vii:vii:subsec:refinamiento-respuesta}

\begin{theorem}[Transport of the differential by refinement]
\label{vii:vii:thm:refinamiento-respuesta}
Let \(J_m:\mathcal H_m\to\mathcal H_{m+1}\) and
\(K_m:\mathcal K_m\to\mathcal K_{m+1}\) be fixed maps.
For a family \(C^1\), suppose the refinement identities
\begin{equation}
 D_{m+1}(t)J_m=K_mD_m(t),\qquad
 K_m^*\mathsf W_{m+1}(t)K_m=\mathsf W_m(t).
 \label{vii:vii:eq:familia-refinamiento-memoria}
\end{equation}
Then \(E_{m+1}(J_mf)=E_m(f)\).
The transport covectors agree on variations linked by
the first identity. When the weights vary, their contributions
to the differential also agree.
\end{theorem}

\begin{proof}
By substitution,
\[
 \langle D_{m+1}J_mf,\mathsf W_{m+1}D_{m+1}J_mf\rangle
 =\langle D_mf,K_m^*\mathsf W_{m+1}K_mD_mf\rangle=E_m(f).
\]
Differentiating the first identity of
\eqref{vii:vii:eq:familia-refinamiento-memoria} gives
\(\dot D_{m+1}J_m=K_m\dot D_m\). Consequently,
\[
 2\operatorname{Re}
 \langle D_{m+1}J_mf,\mathsf W_{m+1}\dot D_{m+1}J_mf\rangle
 =2\operatorname{Re}\langle D_mf,\mathsf W_m\dot D_mf\rangle.
\]
The derivative of the second identity is
\(K_m^*\dot{\mathsf W}_{m+1}K_m=\dot{\mathsf W}_m\), proving
agreement of the weight-variation terms.
\end{proof}

The statement uses compatibility of the transport family.
If \(J_m\) depends on the parameter, the chain rule also
preserves the term
\[
 2\operatorname{Re}
 \langle D_{m+1}J_mf,\mathsf W_{m+1}D_{m+1}\dot J_mf\rangle.
\]
The resulting naturality jointly transports the functional, its
incidences and its differential.

\subsection{Material realization of the response}
\label{vii:vii:subsec:realizacion-respuesta-discreta}

The result of this section is the explicit covector
\(\mathfrak j_m\), with its representative \(\mathcal J_a\).
The material realization composes that covector with the variation of
transport induced by the connection and with the normalization of the action.
If \(\mathcal L_m\) denotes that differential map, the covector on
connection variations is \(\mathcal L_m^*\mathfrak j_m\).
The composition rule
\eqref{vii:vii:eq:variacion-transporte-compuesto} evaluates its part
corresponding to finite products of transports.

In section~\ref{vii:vii:sec:corriente-respuesta-cartan}, the current
\(\sigma_{IJ}\) is a three-form defined by variation of the material
action. Identification between the two realizations preserves the field,
representation, weight, dimensional normalization and passage to the
limit used by that action. These data allow evaluation of the
contraction entering the torsional density. The present
development provides the discrete differential composed in that calculation.
